\documentclass[12pt,letterpaper]{article}
\usepackage[lmargin=1.2in,rmargin=1.2in,tmargin=1in,bmargin=1in]{geometry}
\usepackage[bf,tiny]{titlesec}
	 \usepackage{fontspec,xltxtra,xunicode}
    	\defaultfontfeatures{Mapping=tex-text}
\usepackage[utf8]{inputenc}
%\usepackage[utf8]{fontenc}
\setmainfont{Junicode}

\usepackage{setspace}
\usepackage{gb4e}
\let\eachwordone=\it %italicise the first line of examples
\usepackage[sort]{natbib}
\bibpunct[: ]{(}{)}{,}{a}{}{,}
\usepackage{afterpage}
\usepackage{graphicx}


%%% your packages here; please use sparingly and for things you actually need (not for aesthetics)

\title{Diachronica template}
\author{Claire Bowern \\ Executive Editor, \textit{Diachronica}}
%\date{\today}

\begin{document}

\maketitle

\begin{abstract}
    Here's an abstract of about 150 words
\end{abstract}

\noindent\textbf{Keywords:} templates, fonts, examples, bibliography

\section{Introduction}
We've provided this template for \LaTeX{} users to submit articles to \textit{Diachronica}. Note that it's very minimal. That's deliberate! It is far easier for the typesetters to typeset your article if there is minimal styling.\footnote{Here's an example of a footnote. Note that we do acknowledgements at the end of the paper, not as footnotes to the title or the first section.} There are a number of points about this template which do not conform to the final style for published articles. However, this template is optimized to make the preparation of the submitted manuscript as straightforward as possible.

Here's an example of a sentence with glossed text.

\begin{exe}
\ex Bardi \hfill\citep[123]{bow12}
\gll Ardiyooloon Bardi ngaanka\\
One.Arm.Point Bardi language\\
\glt `One Arm Point Bardi'
\end{exe}

Note that the translation has single quotes. However, a quotation from an individual ``should have double quotes,'' (me, pers. comm. 2021).

Here's an example of a table:
\begin{table}[ht]
    \centering
    \caption{Caption} %note table captions go above the table
    \label{tab:my_label}    \begin{tabular}{cc}
      one   & a \\ \hline
    two     & b
    \end{tabular}
\end{table}

And here's an example of a figure:

\begin{figure}[ht]
    \centering
 %   \includegraphics{}
    \caption{Caption} %note figure captions go below the figure
    \label{fig:my_label}
\end{figure}

\section{Other notes}
This template should use XeLaTeX and unicode: here's ŋ, an engma. Here's [daɪəkʰɹʌnɪkʰə]. Feel free to change the font if you don't like junicode.

\section{Conclusion}
In conclusion, we use the \texttt{unified} stylesheet for linguistics. For further details, please consult the style guide, available from the journal's homepage at \texttt{https://benjamins.com/catalog/dia}.

\section*{Funding information}
Grant or other funding acknowledgements should go here.

\section*{Acknowledgements}
Thanks to Juhyae Kim for comments on this template.

\section*{Abbreviations}
A list of abbreviations used in the text should appear here.


\bibliographystyle{unified}
\bibliography{bibliography}


\bigskip

\appendix
\section{Information about appendices}
If you have any Appendices, this is where they would appear.

\section*{Résumé}

French summary goes here



\section*{Zusammenfassung}

German summary goes here

\bigskip

\section*{Author's address}
\medskip

\noindent\begin{minipage}{\linewidth}
Claire Bowern\\
Department of Linguistics\\
Yale University\\
370 Temple St\\
\textsc{New Haven}, CT\\[3ex]
\texttt{claire.bowern@yale.edu}
\end{minipage}

\end{document}
